\section{Problema e importancia}
En el sistema educativo costarricense, la inclusi\'on de estudiantes con discapacidad se gestiona mediante las Aulas Integradas y los Servicios de Apoyo Educativo. A pesar de que el Ministerio de Educaci\'on P\'ublica (MEP) ha lanzado la ``Gu\'ia Docente: IA en clase (2026)'', el D\'ecimo Informe Estado de la Educaci\'on (2025)\cite{PEN2025} advierte que el pa\'is atraviesa una ``emergencia educativa'' caracterizada por un grave rezago en competencias b\'asicas y una incapacidad del sistema para personalizar el aprendizaje de manera efectiva. 

El problema central es la desconexi\'on entre el uso de herramientas de IA gen\'ericas y la necesidad de una personalizaci\'on profunda que considere el estado emocional del estudiante y la progresi\'on cognitiva según la Taxonom\'ia de Bloom, factores que el informe 2025 identifica como esenciales para superar la crisis de aprendizaje.

Este problema es cr\'itico bajo la \'optica del Estado de la Naci\'on 2025:

\begin{itemize}
    \item Rezago Hist\'orico: El informe 2025 señala que la exclusi\'on y el bajo rendimiento en estudiantes vulnerables han alcanzado niveles alarmantes, poniendo en riesgo la movilidad social. Seg\'un las pruebas PISA 2022, el estudiantado presenta niveles de comprensi\'on lectora y razonamiento matem\'atico que corresponden a niveles de primaria, con un retroceso de hasta diez años \cite{PEN2025}.
    \item Invisibilidad del Estudiante: Existe una incapacidad t\'ecnica para diagnosticar y atender las brechas individuales en tiempo real, lo que el informe denomina ``el apag\'on educativo'' de la personalizaci\'on \cite{MEP2026}.
    \item Debilitamiento institucional: El desmantelamiento de programas consolidados como PRONIE-MEP-FOD (Programa Nacional de Inform\'atica Educativa - Ministerio de Educaci\'on P\'ublica - Fundaci\'on Omar Dengo) ha dejado al sistema sin una pol\'itica clara de inclusi\'on digital y acompañamiento pedag\'ogico en tecnolog\'ia \cite{PEN2025}.
    \item Eficiencia del Gasto: Con la presi\'on fiscal en Costa Rica, optimizar la labor del docente mediante IA constituye una v\'ia viable para escalar la educaci\'on inclusiva sin requerir un aumento en la planilla estatal \cite{MEP2026}.
\end{itemize}